\clearpage

% --------------------------------------------------------------------------------------------

\section {Summary}
\label{sec:summary}

In this technical note,
we have presented the results of a 3-flavour $\nu_{\mu}$-disappearance analysis,
performed by the Valencia-Lancaster-Oxford-Rutherford (VaLOR) group,
on the combined Run1+2+3 (3.010$\times$10$^{20}$ POT) dataset.
This analysis is an update of the work presented in
Refs. \cite{OurRun1DisappAnalysis}, \cite{OurRun12DisappAnalysis} and \cite{OurRun12DisappAnalysis3Flavour}
with improved methods and an enlarged dataset.

This analysis predicts 204.75 $\pm$ 16.75 (syst)
1 $\mu$-like ring events in SuperK in the absence of any oscillation,
but only 58 were observed. The observed deficit has a strong energy dependence;
the ratio of observed to expected, under the no oscillation hypothesis, is
$\sim$30\% $<$ 0.5 GeV, 
$\sim$10\% between 0.5 and 1 GeV and 
$\sim$75\% $>$ 1 GeV.

A $\nu_{\mu}$-disappearance analysis was performed 
in a framework of 3-flavour oscillations including matter effects in constant-density matter.
A total of 48 efficiency, energy scale, neutrino flux, neutrino cross-section,
initial-state nuclear environment and final-state rescattering systematics was taken into account.
The effect of these systematic uncertainties on the expected number of 1 $\mu$-like ring events,
on their reconstructed energy spectrum and on our oscillation measurement were presented in great detail.

Two versions of the oscillation fit were performed: the `2+N' fit allowed all 48 systematics considered 
in this analysis to float, whereas the `2+8' fit allowed only the 8 most dominant systematics to float. 
The dominant systematics are those identified as having an effect on the measured oscillation parameters 
that is greater than 25\% of the expected statistical error, for at least one set of true oscillation 
parameters in the area of interest. The `2+N' 3-flavour $\nu_{\mu}$-disappearance fit to
the observed reconstructed energy spectrum of 1 $\mu$-like ring events of the combined Run1+2+3 dataset,
for $\sin^2\theta_{23}$ yield
$|\Delta m^{2}_{32}|$ = 2.441$\times$10$^{-3}$ $eV^{2}/c^{4}$  and
$\sin^{2}\theta_{23}$ = 0.514
as the best-fit values (p-value = 0.80).
For the '2+N' fit assuming the first octant ($\sin^22\theta_{23} \geq 0.5$),
with $sin^{2}2\theta_{23}$ constrained to the physical region,
the best-fit values are almost identical at
$|\Delta m^{2}_{32}|$ = 2.443$\times$10$^{-3}$ $eV^{2}/c^{4}$  and
$\sin^{2}2\theta_{23}$ = 1.000
(p-value = 0.83).
For the '2+N' fit assuming the second octant ($\sin^22\theta_{23} \geq 0.5$),
the best-fit values are almost identical at
$|\Delta m^{2}_{32}|$ = 2.441$\times$10$^{-3}$ $eV^{2}/c^{4}$  and
$\sin^{2}2\theta_{23}$ = 0.999
(p-value = 0.82).
For the `2+8' fit, the best-fit values are almost identical.
% at 
% $|\Delta m^{2}_{32}|$ = 2.442$\times$10$^{-3}$ $eV^{2}/c^{4}$  and
% $\sin^{2}2\theta_{23}$ = 1.000
% (p-value = 0.82).
All fits used the 2012 PDG values for
$\sin^{2}2\theta_{13}$, $\sin^{2}2\theta_{12}$ and $\Delta m^{2}_{21}$ \cite{PDG2012},
and assumed $\delta_{CP}$=0 and the normal mass hierarchy.

Confidence regions obtained using both the constant-$\Delta\chi^2$ method and the 
method of Feldman and Cousins were presented for both the `2+N' and `2+8' fits.
% Using the marginalized $\Delta\chi^2$ distributions obtained from the fits,
% and the constant-$\Delta\chi^2$ method,
% the 90\% CL allowed region, for both `2+N' and `2+8' fits, assuming the first octant, can be summarized as
% 2.14$\times$10$^{-3}$ $eV^{2}/c^{4}$ $< |\Delta m^{2}_{32}| < $ 2.76$\times$10$^{-3}$ $eV^{2}/c^{4}$ and
% $\sin^{2}2\theta_{23} > $ 0.957.
% The 90\% CL allowed region, for the '2+N' fit, assuming the second octant, can be summarized as
% 2.??$\times$10$^{-3}$ $eV^{2}/c^{4}$ $< |\Delta m^{2}_{32}| < $ 2.??$\times$10$^{-3}$ $eV^{2}/c^{4}$ and
% $\sin^{2}2\theta_{23} > $ 0.???.
Using the results of the '1+N' fits, 
and the constant-$\Delta\chi^2$ method,
the 90\% CL allowed region, for the `2+N' fit for $\sin^2\theta_{23}$, can be summarized as
$2.19$$\times$10$^{-3}$ eV$^{2}$/c$^{4}$ $< |\Delta m^2_{32}| < 2.72$$\times$10$^{-3}$ eV$^{2}$/c$^{4}$ and
$0.397 < \sin^{2}\theta_{23} < 0.630$.
The corresponding limits for the '2+N' fit assuming the first octant, can be summarized as
$2.19$$\times$10$^{-3}$ eV$^{2}$/c$^{4}$ $< |\Delta m^2_{32}| < 2.72$$\times$10$^{-3}$ eV$^{2}$/c$^{4}$ and
$\sin^{2}2\theta_{23}$ $>$ 0.958.
The corresponding limits for the '2+N' fit assuming the second octant, can be summarized as
$2.19$$\times$10$^{-3}$ eV$^{2}$/c$^{4}$ $< |\Delta m^2_{32}| < 2.72$$\times$10$^{-3}$ eV$^{2}$/c$^{4}$ and
$\sin^{2}2\theta_{23}$ $>$ 0.931.

Numerous comparisons were made of the best-fit values and spectra, of the systematic parameter pulls,
and of the $\chi^{2}$ surfaces obtained by the different fit methods and for different running periods.
Detailed comparisons were presented of the confidence regions obtained
by using different construction methods, for the same fit, and
by using different fitting methods, for the same construction method.
The contours from different running periods were compared.
Comparisons of our contours with
recent measurements from MINOS and SuperK and the T2K 3.010$\times$10$^{20}$-POT sensitivity
were also made. Finally, a more critical view of the 4 different results obtained by
this analysis was presented in Sec.\ \ref{sec:primary_result_selection}.



% --------------------------------------------------------------------------------------------

%
%\section* {Acknowledgements}
%\label{sec:acknowledgements}
%
